% Methods: ethics, participants, design, stimuli, procedure, measures, registrations and analyses of
% the five studies, in outline. The detail is in the Supplementary Methods
% (supplement/sections/supp_methods.tex) and in each study's section of the Supplementary Information.
% Every count and result comes from a macro in generated/numbers.tex; design facts (dates, scale
% points, stated hours, viewing seconds, numbers of images, analysis settings) are typed.

\subsection*{Ethics and participants}

All five studies were approved by the University of California, Los Angeles (UCLA) Institutional Review Board (protocol IRB\#24-001292). Each study was an online Qualtrics survey, and participants gave informed consent by choosing `I consent to this study' after reading the consent form.

All participants were students in the UCLA Psychology subject pool and took part for course credit; every analyzed record came through the pool.

The studies are numbered in the order of the paper's argument, not the order in which they ran. Study~2 ran first, then Studies 4, 1 and 3, and Study~5 last. A power analysis calculates how many participants a study needs to detect an effect of a given size. None was run for any study. Each sample size came from a quota, a planned number of participants for each label group (in Study~5, for each combination of label and editing time; Supplementary Methods).

The exclusion rules were written into each registration after the data had been collected. They removed survey test runs, unfinished sessions, anyone who did not answer `Yes, I did' when asked at the end whether they had taken part seriously and, in Studies 2 to 5, anyone who answered an arithmetic question wrongly (the Supplementary Methods give every rule). The rules left \SOneN{} of \SOneNRecords{} exported records in Study~1, \STwoN{} of \STwoNRecords{} in Study~2, \SThreeN{} of \SThreeNRecords{} in Study~3, \SFourN{} of \SFourNRecords{} in Study~4 and \SFiveN{} of \SFiveNRecords{} in Study~5. These \NTotal{} participations came from \NPeople{} people. Of these, \NPeopleRepeat{} took part twice, each in one earlier study and then in Study~5, so they had already seen some of the same originals and read the debriefing, which told them that anything produced with AI is insufficiently creative (Procedure). No registered analysis removes them, no analysis checks whether their first session changed their ratings, and the exploratory model that pools Studies 2, 4 and 5 counts twice those who had taken part in Study~2 or~4.

The mean ages in Studies 1 to 5 were \SOneAgeMean, \STwoAgeMean, \SThreeAgeMean, \SFourAgeMean{} and \SFiveAgeMean{} years, and \SOnePctFemale, \STwoPctFemale, \SThreePctFemale, \SFourPctFemale{} and \SFivePctFemale{} of participants identified as female; the Supplementary Methods give the full counts.

\subsection*{Design and stimuli}

The label told participants how the artist, Alex, had made each original image. Under the Human label he used his graphic design skills and digital sketching tools. Under the AI label he developed and input prompts into an AI system, which generated the artworks (the Supplementary Methods quote each version). The survey assigned each participant at random, each label equally often, to one label only. Both label groups saw the same images, presented in the same way, so a difference between the groups is an effect of what they were told. The Supplementary Methods give the size of each label group.

Studies 1 and 3 asked about the originals. Study~1's participants rated all ten. Study~3's rated nine, chosen at random, three at each of three stated times that Alex had spent on them: under 1 hour, 10 hours or 100 hours. Under the Human label he spent the hours drawing, and under the AI label revising prompts.

Studies 2, 4 and 5 asked about second works. Both label groups read, in the same words, that Alex had made each one from an original with his design and sketching tools. Each original was followed by three second works, each shown beside it: None, the original shown again unchanged; Slight, a light edit; and Dramatic, a heavy edit. Study~2 used all ten originals. Studies 4 and 5 used six, chosen at random, three with a human or humanoid figure and three without. Study~4 also stated how long Alex had spent on each original, which we call source time, with two originals at each of the three times. Study~5 stated these times too, but a fault in the survey script left their record blank, so no result about source time can be reported for Study~5. Study~5 also stated how long Alex had spent making the second works, which we call editing time. Each participant was told one editing time for the whole session, assigned at random with the label, so the six combinations formed separate groups of \SFiveNCellMin{} to \SFiveNCellMax{} people. Stated time, source time and modification level varied within participants, each of whom saw every level.

The ten originals were made with OpenAI's DALL\textperiodcentered E image generator, used through ChatGPT, from prompts that we wrote; no project file records their wording. Eight of the ten image files carry a provenance record, a statement of origin written into the file and digitally signed by OpenAI under the standard of the Coalition for Content Provenance and Authenticity (C2PA). Each of the eight names DALL\textperiodcentered E and ChatGPT, marks the image as AI-generated and is dated 6 to 8 November 2024. The other two files carry no record. Five originals show a human or humanoid figure and five do not; the analyses call this the kind of image. Every image under both labels was made with AI, so the Human label was a cover story.

We made the Slight and Dramatic versions by editing each original by hand, in image-editing software that we operated manually: Adobe Photoshop and, for one Dramatic version, the online editor Dzine, used only as a manual editor, whose logo we then removed from that image. No generative or other AI editing tool changed any image. The registrations of Studies 2, 4 and 5 describe the Slight versions as color and contrast adjustments, but \XSImageCheckSlightReframedN{} of the ten were also cropped or resized (Fig.~\ref{fig:design}a; Supplementary Methods).

\subsection*{Procedure}

Before every original, a short description restated the label; in Studies 3 to 5 it also stated a time (Supplementary Methods). The original then stayed on screen alone for at least 8 seconds. In Studies 2, 4 and 5, each second work was then shown beside it for at least 10 seconds. The ratings came on the next screen, without the image, in a fixed order, with a required box below them for explaining each rating. Every explanation was therefore written after its rating, although a rating could still be changed before the screen was left. The survey did not allow going back.

The originals and the second works came in random orders, not all of which were recorded, so no analysis takes order into account. No manipulation check asked whether participants had noticed, remembered or believed the label.

The final screen, a debriefing that was the same in all five studies, said: `Current legal standards state that anything produced with artificial intelligence is insufficiently creative, even if you use AI as a base against which you create new work.' That overstates the law: the US Copyright Office registers a person's own sufficiently creative selection, arrangement or modification of AI-generated material, with that material left out of the claim.\cite{usco2023guidance,usco2025copyrightability} The debriefing thus stated as law the categorical view that this paper calls the original sin and tests as a hypothesis about lay judgment, and it did not tell participants that every image had been made with AI (Supplementary Methods).

\subsection*{Measures}

The ratings are lay judgments of creativity, effort, transformativeness and essence change. None is a judgment of authorship, of originality in the legal sense, or of whether a work can be protected. Each item offered five labeled answers, scored 1 to 5 by position. Originals were rated for creativity and for effort (the Results quote both items). Only the creativity item carried a definition, and neither says whose creativity or effort is meant. The creativity definition echoes copyright's standard of originality, at least a minimal degree of creativity (\emph{Feist Publications, Inc.\ v.\ Rural Telephone Service Co.}, 499 U.S.~340, 345 (1991)). That standard is a floor, not a scale, so a rating of how much creativity was involved is not a judgment of whether a work clears it.

Second works were rated for transformativeness, essence change and creativity. The transformativeness item defined transformativeness as `the degree the second work adds new expression, meaning, or message to the first work'. Studies 4 and 5 also defined essence and creativity, and Study~2 defined neither (Supplementary Methods). Effort was not rated for second works. The transformativeness definition borrows the Supreme Court's description of a transformative use in fair-use law (\emph{Campbell v.\ Acuff-Rose Music, Inc.}, 510 U.S.~569, 579 (1994)). Here it is a lay rating of what the second work adds, not a fair-use judgment. In each of Studies 2, 4 and 5, two transformativeness items had a sixth answer point. As registered, their ratings (\STwoNSixPointDropped, \SFourNSixPointDropped{} and \SFiveNSixPointDropped) are left out of the primary analyses.

\subsection*{Registrations}

All five studies were registered on the Open Science Framework (OSF), an online archive for study plans, on 8 October 2026 (\url{https://osf.io/u3d69/overview?view_only=be795b275c3548428d74457df31a4159}), after their data had been collected, so no study was preregistered. What we had already seen of each study's data before filing, and the corrections to our earlier report of Studies 2 and 4, are set out in the Supplementary Information (Registrations).

\subsection*{Statistical analysis}

The main quantity is the label's penalty: the mean rating under the Human label minus the mean under the AI label, in points on the 1-to-5 scale. No condition lacked a label, so a penalty cannot be split into a drop under one label and a rise under the other.

The registered analyses fitted a linear mixed model to every individual rating of each question, leaving missing ratings out (the Supplementary Methods give the counts). Each model gave every participant and every original its own baseline, called a random intercept. A participant or an original rated high throughout therefore did not distort comparisons between conditions. Each model included the factors that its study varied and all their interactions, an interaction being a change in one factor's effect across the levels of another. The kind of image was crossed with those factors in Studies 1 and 2, added without interactions in Study~5 and left out in Studies 3 and 4 (the Supplementary Methods give each model). The models did not let the label's effect differ from one original to the next. Their intervals therefore describe these ten images and do not allow for how the penalty might vary across other images.

Each reported difference compares the model's predicted mean ratings for two conditions, each averaged with equal weight over the other factors. Tests were two-sided, counting a difference in either direction, at .05. Their degrees of freedom, which set how large a statistic must be to count as significant, came from Satterthwaite's approximation. Study~1's contrasts used the Kenward--Roger approximation, which gives the same degrees of freedom for its penalty (Supplementary Methods).

Within each study, the tests that its registration names as confirmatory were corrected together by Holm's method.\cite{holm1979simple} The corrected sets held \SOneNConfirmatory, \STwoNConfirmatory, \SThreeNConfirmatory, \SFourNConfirmatory{} and \SFiveNConfirmatory{} tests in Studies 1 to 5. Study~1's four are H1 on each rating, H2 and H4. H4, whether the AI label spreads ratings more between participants and less within each participant, is decided by two ratios of variances, a measure of spread (Supplementary Methods). It enters the corrected set as a likelihood-ratio test of whether letting the variance between participants differ by label improves the model's fit. Study~1's H3, whether the penalty is the same for originals with and without a figure, is decided by its own registered rule outside the set. So are Study~3's confirmatory equivalence tests and decision rule.

Where a registration asked whether an effect was small, it used an equivalence test, two one-sided tests at .05 that pass when the effect's 90\% confidence interval lies wholly inside a stated bound.\cite{schuirmann1987comparison,lakens2017equivalence} In Studies 2 and 4, the larger one-sided $p$-value also had to survive Holm's correction.

Three bounds were used, in rating points. Study~1's H3 used $\pm$\SOneHThreeBound, a fifth of a scale point, set after its analyses had been run. Studies 2, 4 and 5 used $\pm$\SourceBound{} for the label's effect on all three second-work measures, and Study~3 used it for the label's gap at 100 hours. It is half of Study~1's registered creativity penalty of \SOneHOneCreatDiffPrecise{} points. The bound is a fixed number, although the penalty it halves was itself estimated with uncertainty (95\% CI \SOneHOneCreatLo{} to \SOneHOneCreatHi). Each registration set it knowing Study~1's estimate and, for Studies 2 to 4, the earlier analyses of its own data. Study~5 tested editing time against $\pm$\SFiveEditTimeBound, Study~3's registered creativity contrast between 100 hours and under 1 hour, by a rule fixed before any Study~5 model was fitted (Supplementary Methods).

\subsection*{Model checks and further registered analyses}

As registered checks, the confirmatory models were refitted as ordinal models (cumulative-link mixed models), which treat the answers as ordered categories. In Study~1, agreement of this ordinal model was part of the registered rule for H2. Elsewhere a disagreement is reported beside the verdict rather than replacing it. A model that did not converge yields no $p$-value, interval or verdict, and its check is reported as not completed; every study's code applies this rule. A singular fit keeps its numbers and is reported as singular. Supplementary Note~2 and the Supplementary Methods describe the other registered analyses, among them checks in models that let effects vary between participants and between originals, Study~1's robustness checks, refits of Studies 1 and 4 without the Human-label participants who wrote that a work was machine-made, and Study~5's rescaled reanalysis.

\subsection*{The written explanations}

Eight fixed word lists, sets of words and phrases that a script searches for in every explanation in the same way, flagged what each explanation mentioned, such as visible detail, time, skill, prompting or machines, or a statement that a machine made the work (the Supplementary Methods list all eight). A flag records that a topic was mentioned, not that it was endorsed. We wrote the lists by reading Study~1's explanations with the label and the rating hidden. Four parts of them, among them the whole list for a statement that a machine made the work, were settled after results by label had been read. Studies 2 to 5 ran the lists unchanged. The script received only identifiers and text. No person and no language model labeled, coded or judged any explanation for these measures. An earlier language-model coding of Study~1's explanations, made with the label in its instructions, was retired before registration and is not used.

Each registration names two registered secondary checks on the explanations. All other analyses of them are exploratory, among them Study~1's five tests corrected by the Benjamini--Hochberg method,\cite{benjamini1995controlling} whose results are descriptive readings, not verdicts (Supplementary Methods). Because each explanation was written after its rating, none of these measures can show why anyone rated as they did.

\subsection*{Exploratory analyses}

The exploratory analyses were specified on 9 and 10 October 2026, after every registered result was known. Their tests are not corrected as a set, and none changes a registered verdict. The Supplementary Methods and Supplementary Notes~4 to~7 describe them.

\subsection*{Software and reproducibility}

The models ran in R 4.5.2 (lme4 2.0-1, lmerTest 3.2-1, emmeans 2.0.3, ordinal 2025.12-29 and brms 2.23.0), and the text measures and the image check in Python 3.12.3. The Supplementary Methods give the other versions, those of the pretrained models and the fixed starting values of the random-number generators. Each analysis records a checksum, a fingerprint of a file's exact bytes, for every input and output, and a script reads every result, and every count of participants, records or ratings, from the analysis tables and records the table, row and column behind each.

% The authors will add the statement on the use of large language models here (journal requirement).

\subsection*{Reporting summary}

Further information on research design is available in the Nature Portfolio Reporting Summary linked to this article.
